% =============================================================================
% Lecture 13: Workload Placement
% Physical AI: Machine Learning Systems That Sense and Act (Volume IV), Chapter 13
% =============================================================================
% Thesis: separating the proposer and the permission path in software does not
% make them independent while they share memory, power, clocks, resets, or
% heat, so a placement is a hypothesis that loaded tests on the real silicon
% must try to refute.
% Running case: the warehouse mobile manipulator's 1 kHz permission tick
% (unloaded 135 us, declared WCET 250 us, deadline 400 us) on a discrete
% lockstep MCU, against the same tick placed hypothetically on the application
% SoC. Hook and synthesis: the thermal hangover (Napkin Math 13.2).
%
% Build: make deck V=4 D=13_placement   (figures: images/make_figs.py)
% Plan:  lecture_map.yml in this folder. Numbers: images/calcs.py.
% =============================================================================

\documentclass[aspectratio=169, 12pt]{beamer}
\usepackage{../../assets/beamerthememlsys}
\mlsyssetup{
  volume={Volume IV},
  chapter={Chapter 13},
  chaptertitle={Workload Placement},
  layer=nervous,
  logo={../../assets/img/logo-mlsysbook.png},
  instlogo={../../assets/img/logo-harvard.png},
}
\usepackage[T1]{fontenc}
\usepackage[scaled=0.9]{helvet}
\usepackage{courier}
\renewcommand{\familydefault}{\sfdefault}
\usepackage{booktabs,amsmath,array}
\graphicspath{{images/}}
% Generated by calcs.py. Do not edit; edit calcs.py and rerun it.
\defval{tick_khz}{1}
\defval{t_tick_ms}{1}
\defval{mcu_mhz}{400}
\defval{chunk_infer_ms}{40}
\defval{chunk_period_ms}{50}
\defval{sense_us}{40}
\defval{check_us}{80}
\defval{stage_us}{15}
\defval{unloaded_us}{135}
\defval{unloaded_cycles}{54{,}000}
\defval{wcet}{250}
\defval{deadline}{400}
\defval{blocking}{50}
\defval{interference}{40}
\defval{b_plus_i}{90}
\defval{exec_ceiling}{310}
\defval{allowance_us}{60}
\defval{deadline_share}{0.4}
\defval{rail_holdup_s}{2}
\defval{backlog_bytes}{360{,}448}
\defval{backlog_tx}{88}
\defval{tile_tx}{64}
\defval{tile_kib}{256}
\defval{snapshot_kib}{4}
\defval{bw_gbs}{204}
\defval{bank_switches}{3}
\defval{conflict_ns}{37.5}
\defval{t_rcd_min_ns}{18}
\defval{t_ck_ns}{1.25}
\defval{rcd_cycles}{15}
\defval{data_rate}{6{,}400}
\defval{backlog_us}{1.879}
\defval{ladder_backlog_us}{1.88}
\defval{compute_after_us}{95}
\defval{n_sense}{21}
\defval{read_room_wcet}{155}
\defval{n_wcet}{82}
\defval{read_room_dead}{215}
\defval{n_dead}{114}
\defval{end_dead}{309}
\defval{channels}{16}
\defval{backlog1_us}{28.4}
\defval{n1_wcet}{5}
\defval{valsan_a9}{2.1}
\defval{valsan_neh}{7.9}
\defval{valsan_a15}{10.6}
\defval{valsan_a15_w}{21.4}
\defval{bechtel_max}{346}
\defval{stretch_wcet}{1.85}
\defval{stretch_wcet_less1}{0.85}
\defval{stretch_dead}{2.30}
\defval{s_mshr}{21}
\defval{phi21}{4.3}
\defval{phi21_dead}{6.5}
\defval{s_wb}{300}
\defval{phi300}{0.28}
\defval{s_ten}{10}
\defval{s_ten_less1}{9}
\defval{phi10}{9.5}
\defval{phi10_us}{12.8}
\defval{orin_poll_ms}{1{,}000}
\defval{t_amb}{45}
\defval{theta}{2}
\defval{tau_s}{8}
\defval{p_quiet}{4}
\defval{p_load}{26}
\defval{t_quiet}{53}
\defval{t_loaded}{97}
\defval{t_trip_c}{95}
\defval{t_clear_c}{85}
\defval{t_trip_s}{24.7}
\defval{trip_head}{42}
\defval{clear_head}{32}
\defval{recover_s}{2.18}
\defval{recover_ticks}{2{,}175}
\defval{recover_loaded_s}{2.55}
\defval{f_nom_ghz}{1.2}
\defval{f_thr_mhz}{600}
\defval{tick_cycles}{300{,}000}
\defval{thr_us}{500}
\defval{resp_nom}{340}
\defval{resp_thr}{590}
\defval{past_us}{190}
\defval{overrun_pct}{47.5}
\defval{theta_max}{1.92}
\defval{di_a}{6}
\defval{rise_ns}{2}
\defval{l_ph}{30}
\defval{r_mohm}{8}
\defval{v_nom}{0.85}
\defval{droop_l_mv}{90}
\defval{droop_r_mv}{48}
\defval{droop_mv}{138}
\defval{delay_pct}{39}
\defval{lease_ms}{60}
\defval{min_rate_hz}{16.7}
\defval{bechtel_unreg}{345}
\defval{bechtel_w100}{1.34}
\defval{bechtel_w100_1}{1.3}
\defval{bechtel_w50}{1.2}
\defval{bao_linux_cpus}{3}
\defval{bao_linux_gib}{1}
\defval{bao_rtos_cpus}{1}
\defval{bao_rtos_mib}{128}
\defval{bao_shmem_kib}{64}
\defval{chunk_bytes}{4{,}096}
\defval{spi_mhz}{25}
\defval{link_ms}{1.31}
\defval{deser_ms}{2}
\defval{valid_ms}{18}
\defval{activate_ms}{1}
\defval{d2a_ms}{63.3}
\defval{d2a_mm}{82.3}
\defval{v_aisle}{1.3}
\defval{spare_mm}{131.0}
\defval{ceil_ms}{100.8}
\defval{reserve_ms}{37.5}
\defval{reserve_mm}{48.7}
\defval{link_share}{2}
\defval{plus_infer_ms}{103.3}
\defval{small_ms}{4}
\defval{jitter_ms}{5}
\defval{stretch_pct}{12.5}
\defval{renew_window_ms}{55}
\defval{rt2_5b_hz}{5}
\defval{rt2_5b_ms}{200}
\defval{rt2_ratio}{3.6}
\defval{rt2_55b_low}{1}
\defval{rt2_55b_high}{3}
\defval{fogros_local_s}{14}
\defval{fogros_off_s}{1.2}
\defval{wd_ms}{5}
\defval{wd_mm}{6.5}
\defval{ex_tick}{260}
\defval{ex_before_deadline}{140}
\defval{tick_mm}{1.3}
\defval{cp_us}{300}
\defval{cp_over}{50}
\defval{cp_resp}{390}
\defval{cp_left}{10}
\defval{trace_exp}{7}
\defval{trace_millions}{10}
\defval{trace_h}{2}
\defval{trace_min}{46}
\defval{trace_sec}{40}
\defval{cold_p50}{95}
\defval{cold_p99}{128}
\defval{cold_max}{152}
\defval{hot_p50}{210}
\defval{hot_p99}{330}
\defval{hot_p999}{395}
\defval{hot_max}{468}
\defval{tick_ns}{1{,}000{,}000}
\defval{deadline_ns}{400{,}000}
\defval{renewal_onset_ms}{82}

\setcounter{mlsystotalsections}{5}

\title{Workload Placement}
\author{Vijay Janapa Reddi}
\institute{Harvard University}
\date{}

% Volume IV palette colors for text beside figures
\definecolor{vnavy}{HTML}{1F407A}
\definecolor{vgreen}{HTML}{047857}
\definecolor{vred}{HTML}{991B1B}
\definecolor{vamber}{HTML}{B45309}
% Notation used on the slides
\newcommand{\Twcet}{T_{\text{WCET}}}
\newcommand{\Tdead}{T_{\text{deadline}}}
\newcommand{\Tload}{T_{\text{loaded}}}
\newcommand{\Tiso}{T_{\text{isolated}}}
\newcommand{\Dint}{\Delta_{\text{shared}}}
% A large answer choice for predict and concept-question slides
\newcommand{\choice}[2]{\textbf{\textcolor{mlsysaccent}{#1}}\enspace #2}
\newcommand{\answerat}[3]{\alt<#1->{\colorbox{mlsysaccent!15}{\choice{#2}{#3}}}{\choice{#2}{#3}}}
% A goal or takeaway line with its thumbnail
\newcommand{\thumbline}[2]{%
  \begin{minipage}[c]{2.6cm}\centering\includegraphics[width=2.5cm,height=0.95cm,keepaspectratio]{#1}\end{minipage}%
  \hspace{0.3cm}\begin{minipage}[c]{0.74\textwidth}\small #2\end{minipage}}
% An exact excerpt of real source, at a legible size
\newenvironment{logblock}[1][7.5]{\par\ttfamily\fontsize{#1}{9}\selectfont\raggedright\setlength{\parskip}{0pt}}{\par}
\newcommand{\ind}[1]{\hspace*{#1\fontcharwd\font`0}}

\begin{document}

% =============================================================================
% TITLE AND HOOK
% =============================================================================
\mlsysnote{[0.5 min] [CORE] Read the thesis aloud once. This lecture closes Part III, Running the Machine, and its law: proposal is not permission. The last lecture gave the permission path its job on every tick; today asks where that job runs, and whether anything the proposer does can make it late. The notes total 71.5 minutes, 58 of them core. FLEX: cutting every [OPTIONAL] frame leaves a complete lecture of about 58 minutes; cut first the timescale ladder, the configuration file, and the open questions.}
\mlsystitle{Workload Placement}{Where each computation runs decides whether the permission path can still refuse in time.}{cover_placement_blueprint.png}

\begin{frame}{The safety check has its own core. Can the network still make it late?}
  \centering
  \safeimg[width=0.86\textwidth,height=4.9cm,keepaspectratio]{tesla-fsd-board.jpg}\\[0.1cm]
  {\small Two systems-on-chip (SoCs), each with its own memory, on one board supply.}
  \mlsyscredit{Tesla FSD computer board. Photo: Votpuske, CC BY 4.0, Wikimedia Commons.}
  \note{[1.5 min] [CORE] LINK: the last lecture built a permission path that re-checks every command on a \val{tick_khz} kHz tick. NARRATE: this is a real driving computer, Tesla's FSD board. Point to the two large packages, each an SoC with its own memory chips beside it, and to the connectors and power parts that serve both. Our warehouse robot has the same split between a learned proposer that runs hard and a permission check that must answer every millisecond. ENGAGE: ask the class, ``If the safety check runs on its own core, can the neural network still make it miss its deadline?'' Take a show of hands. Most students say no, because the cores and the software are separate. WARN: students equate a separate core with a separate failure domain; the next slide shows a coupling no software diagram contains. FLEX: [CORE].}
\end{frame}

\begin{frame}{Predict: inference stops. When is the tick on time again?}
  \mlsysactivity{Predict}{2 min}
  \centering
  \mlsyssteps[width=0.92\textwidth,height=5.05cm,keepaspectratio]{hook_step}{3}\\[0.0cm]
  {\normalsize\choice{A}{at once}\qquad\choice{B}{next tick}\qquad\answerat{3}{C}{about 2 s later}\qquad\choice{D}{never}}
  \mlsyscredit{Workload Placement, Napkin Math 13.2 (illustrative: \val{t_amb}\,°C bay, \val{theta} K/W, \val{tau_s} s, \val{p_load} W loaded, trip \val{t_trip_c}\,°C, clear \val{t_clear_c}\,°C). Tick on a cluster that shares the SoC's clock governor.}
  \note{[3 min] [CORE] NARRATE: the permission tick placed on a real-time cluster of the application die, which shares the die's clock governor. Step 1: the proposer runs continuous inference; junction temperature rises from \val{t_quiet} toward \val{t_loaded}\,°C with an \val{tau_s} s time constant; every tick ends at \val{resp_nom} us, inside its \val{deadline} us deadline. Step 2: at \val{t_trip_s} s the junction reaches the \val{t_trip_c}\,°C trip and the governor halves the clock from \val{f_nom_ghz} GHz to \val{f_thr_mhz} MHz. The same \val{tick_cycles} cycles now take \val{thr_us} us, so the tick ends at \val{resp_thr} us, \val{past_us} us past the deadline, \val{overrun_pct} percent late. ENGAGE: suppose the proposer stops right at the trip. When is the tick on time again? Vote. Step 3: C. The clock returns only when the junction cools to the \val{t_clear_c}\,°C clear point: \val{recover_s} s, \val{recover_ticks} ticks. Students who chose A think of the governor as following the load; it follows the temperature. WARN: these parameters are illustrative and describe no measured part; the shape is the point. If inference keeps running to the \val{t_loaded}\,°C steady state, recovery takes \val{recover_loaded_s} s (derived). The synthesis returns here. FLEX: [CORE].}
\end{frame}

\begin{frame}{Recall: what did the enforcement record promise?}
  \mlsysactivity{Retrieval}{1.5 min}
  \begin{columns}[c]
    \begin{column}{0.56\textwidth}
      {\Large\bfseries Close your notes.}\\[0.2cm]
      {\normalsize What execution time did the permission tick declare, inside what deadline? Where did the volume first place the permission path?}\\[0.25cm]
      \uncover<2->{\small\textbf{\textcolor{mlsysaccent}{Declared worst-case execution time (WCET)}} \val{wcet}\,µs, deadline \val{deadline}\,µs, tick \val{t_tick_ms}\,ms\\[0.05cm]
      \textbf{\textcolor{mlsysaccent}{A discrete safety microcontroller (MCU)}}, apart from the proposer's SoC}
    \end{column}
    \begin{column}{0.40\textwidth}
      \centering\safeimg[width=\textwidth,height=1.5cm,keepaspectratio]{thumb_budget_step3.pdf}\\[0.15cm]
      \uncover<3->{\begin{mlsyscard}{mlsysaccent}\small\textbf{From Lectures 4 and 12:} two processors were a \emph{hypothesis}; the WCET bounds the tick \emph{alone}.\end{mlsyscard}}
    \end{column}
  \end{columns}
  \mlsyscredit{Safety Enforcement, the enforcement record; The Nervous System, the two-processor hypothesis.}
  \note{[1.5 min] [CORE] NARRATE: give 45 seconds. Step 2: Lecture 12's enforcement record declares a worst-case execution time (WCET) of \val{wcet} us inside a chosen \val{deadline} us deadline in a \val{t_tick_ms} ms tick. Lecture 4 placed the permission path on a discrete safety microcontroller (MCU) as a working hypothesis. Step 3: both promises were made about the tick running alone, and the placement was never tested under the proposer's load. ENGAGE: what would make a bound measured alone wrong under load? Take two answers; the lecture answers it. FLEX: [CORE].}
\end{frame}

\begin{frame}{Goals: test a placement against the silicon it shares}
  \thumbline{thumb_board_step4.pdf}{List what two paths still share on a board, row by row, for four candidate placements.}\\[0.2cm]
  \thumbline{thumb_backlog_step5.pdf}{Price memory queueing, shared-cache miss handling, rail droop, and heat against the \val{allowance_us}\,µs allowance.}\\[0.2cm]
  \thumbline{thumb_boardlink_step2.pdf}{Choose a remedy for a measured coupling, and fit the board link inside the stopping budget.}\\[0.2cm]
  \thumbline{thumb_trace_step2.pdf}{Name the measurement that refutes a placement, and what the placement record must carry.}\\[0.15cm]
  {\footnotesize\textcolor{midgray}{Thread: the third law, proposal is not permission. Record: the placement record.}}
  \note{[1 min] [CORE] NARRATE: four goals, one per part of the lecture; each thumbnail is the shape of the figure that teaches it. Part III's law, proposal is not permission, which Lecture 8 stated, separates the two paths in authority (Lecture 1 drew that proposal boundary). Today asks whether the hardware keeps them separate in time when the proposer works hardest. The lecture adds one record to the chain, the placement record. FLEX: [CORE].}
\end{frame}

% =============================================================================
% SECTION: INVENTORY
% =============================================================================
\mlsysnote{[0.5 min] [CORE] First segment: what do two paths on separate cores still share? Lecture 4 owns the independence requirement; this segment argues it resource by resource on real silicon.}
\mlsysdivider{What do two paths on separate cores still share?}{Workload Placement, \S13.1 and \S13.2}
\section{Inventory}

\begin{frame}{Production boards put the safety MCU in its own package}
  \begin{columns}[c]
    \begin{column}{0.68\textwidth}
      \centering\safeimg[width=\textwidth,height=4.9cm,keepaspectratio]{tesla-hw25-board.jpg}
    \end{column}
    \begin{column}{0.30\textwidth}
      \centering\safeimg[width=0.82\textwidth,height=2.3cm,keepaspectratio]{tesla-hw25-tricore.jpg}\\[0.1cm]
      {\footnotesize Infineon AURIX TC297: lockstep cores, two cores that compare every step, sold for the top automotive integrity grade (ASIL-D)}
    \end{column}
  \end{columns}
  \vspace{0.1cm}
  {\small Tesla Autopilot HW2.5: NVIDIA processors propose; a separate microcontroller sits beside them.}
  \mlsyscredit{Photo: Steve Jurvetson, CC BY 2.0, Wikimedia Commons (cropped). Lockstep: Infineon AURIX TC29xT product page.}
  \note{[2 min] [OPTIONAL] LINK: Lecture 4 placed our permission path on a safety microcontroller; industry does the same. NARRATE: Tesla's Autopilot HW2.5 board. The three large packages are NVIDIA processors; the center one is a GP106 GPU. At the upper right, the inset enlarges an Infineon TriCore, part number TC297, from the AURIX family; the silkscreen beside it reads AURIX debug LEDs. Infineon sells it with lockstep cores, in which two cores run the same instructions a few cycles apart and a comparator flags any difference, which catches a fault inside the core. ENGAGE: what does this MCU still share with the NVIDIA chips? Expected: the board, its power input, connectors, and the enclosure's heat path. WARN: the photo does not tell us what software ran on the TriCore; say only that the board carries a separate safety-class MCU. FLEX: [OPTIONAL].}
\end{frame}

\begin{frame}{The permission tick has \val{allowance_us}\,µs to spare for shared hardware}
  \centering
  \mlsyssteps[width=0.94\textwidth,height=4.8cm,keepaspectratio]{budget_step}{4}\\[0.15cm]
  {\small \val{mcu_mhz} MHz lockstep MCU, \val{tick_khz} kHz tick.\enspace Proposer: \val{chunk_infer_ms}\,ms of each \val{chunk_period_ms}\,ms chunk period on the SoC.}
  \mlsyscredit{Workload Placement, \S13.2: sub-costs \val{sense_us} + \val{check_us} + \val{stage_us}\,µs (illustrative); WCET and deadline chosen; \val{blocking} + \val{interference}\,µs blocking and interference (illustrative).}
  \note{[3 min] [CORE] NARRATE: one tick drawn to scale. Step 1: unloaded, the tick samples encoders and the IMU (\val{sense_us} us), checks the stopping envelope (\val{check_us} us), and stages setpoints for the drives (\val{stage_us} us): \val{unloaded_us} us, or \val{unloaded_cycles} cycles at \val{mcu_mhz} MHz. Step 2: the declared worst-case execution time, \val{wcet} us, the bound the enforcement record carries. Step 3: the deadline is \val{deadline} us, but the MCU's own tasks already claim \val{b_plus_i} us of it: \val{blocking} us of blocking by lower-priority work and \val{interference} us of interference from interrupts, the response bound of Lecture 4. So execution may take at most \val{exec_ceiling} us. Step 4: the declared bound sits at \val{wcet}, so shared hardware may add at most \val{allowance_us} us before the deadline breaks. ENGAGE: why charge interference to the tick's execution time rather than as a separate term? Expected: shared hardware slows the tick's own instructions. WARN: the deadline is \val{deadline_share} of the period, so the tick must finish by \val{deadline} us; finishing before \val{t_tick_ms} ms misses it. FLEX: [CORE].}
\end{frame}

\begin{frame}{Separate cores still share the cache, clock, rail, and heat}
  \begin{columns}[c]
    \begin{column}{0.68\textwidth}
      \centering\mlsyssteps[width=\textwidth,height=6.0cm,keepaspectratio]{board_step}{4}
    \end{column}
    \begin{column}{0.30\textwidth}
      \begin{mlsyscard}{mlsysaccent}\footnotesize
        \textbf{Operational independence:} no proposer load or fault can push the tick past its deadline or disable its refusal, through \emph{any} shared resource.
      \end{mlsyscard}
      \vspace{0.1cm}
      {\footnotesize Red: shared. Clock governor: lowers the clock when the die is hot.}
    \end{column}
  \end{columns}
  \mlsyscredit{Redrawn from Workload Placement, \S13.2.2, figure 1.}
  \note{[3.5 min] [CORE] NARRATE: the machine's compute board, built in four steps; each badge is a candidate home for the permission tick. Step 1: (a) shared application cores under Linux beside the neural accelerator. Everything red is shared: last-level cache, DRAM controller, clock governor, core rail, reset chain, fabric, and the package's heat path. Step 2: (b) a real-time cluster has its own cores, but every red block is still shared. Step 3: (c) an on-die safety island with its own crystal, rail, and SRAM behind a firewall; it still shares the die temperature, the fabric that carries proposals, the board supply, and possibly the reset chain. Step 4: (d) a discrete MCU in its own package, with private memory, its own clock, a reset supervisor, an isolated rail with \val{rail_holdup_s} s of hold-up, and a board link. Only the battery remains shared. Read the definition card. ENGAGE: which candidate does memory protection make independent? Expected: none; memory protection governs where an access may go and says nothing about when it is served. WARN: lockstep cores on a shared rail share a failure. FLEX: [CORE].}
\end{frame}

\begin{frame}{Each candidate leaves different rows shared}
  \begin{columns}[c]
    \begin{column}{0.66\textwidth}
      \footnotesize
      \begin{tabular}{@{}l>{\centering\arraybackslash}p{0.95cm}>{\centering\arraybackslash}p{0.95cm}>{\centering\arraybackslash}p{1.4cm}>{\centering\arraybackslash}p{1.4cm}@{}}\toprule
        \textbf{Resource} & \textbf{(a)} & \textbf{(b)} & \textbf{(c)} & \textbf{(d)}\\\midrule
        Cache, miss handling & \textcolor{vred}{shared} & \textcolor{vred}{shared} & separate & separate\\
        DRAM controller & \textcolor{vred}{shared} & \textcolor{vred}{shared} & separate* & separate\\
        Core rail, clock & \textcolor{vred}{shared} & \textcolor{vred}{shared} & separate & separate\\
        Reset chain & \textcolor{vred}{shared} & \textcolor{vred}{shared} & \textcolor{vamber}{evidence} & separate*\\
        Fabric, die heat & \textcolor{vred}{shared} & \textcolor{vred}{shared} & \textcolor{vred}{shared} & link\\
        Board supply & \textcolor{vred}{shared} & \textcolor{vred}{shared} & \textcolor{vred}{shared} & \textcolor{vred}{shared}\\\bottomrule
      \end{tabular}\\[0.05cm]
      {\scriptsize * if the tick stays in island memory; if the supervisor is independent}\\[0.08cm]
      \uncover<2->{\footnotesize\textbf{Jetson AGX Orin module:} island present, \textcolor{vred}{not available} to integrators. \textbf{DRIVE AGX Orin:} available, 3\,MB SRAM.}
    \end{column}
    \begin{column}{0.32\textwidth}
      \centering\safeimg[width=\textwidth,height=3.6cm,keepaspectratio]{jetson-agx-orin-devkit.jpg}\\[0.05cm]
      {\footnotesize Jetson AGX Orin developer kit}
    \end{column}
  \end{columns}
  \mlsyscredit{Workload Placement, \S13.2.2, tables 1 and 2 (NVIDIA 2023, 2026). Photo: Auledas, CC BY 4.0, Wikimedia Commons.}
  \note{[2.5 min] [CORE] NARRATE: the inventory from the board, row by row. Step 1: (a) and (b) share every row; (c) separates cache, DRAM, rail, and clock but keeps the fabric, die temperature, and board supply, and its reset chain needs evidence; (d) leaves only the battery and adds a board link. On the inventory alone, (c) is a strong contender. Step 2: the machine rejects it for a reason outside silicon. The Jetson AGX Orin module, the class of computer in the photo, carries a Cortex-R52 safety island, but NVIDIA does not make it available to integrators on that module; the DRIVE AGX Orin platform does, with 3 MB of island SRAM. An island nobody can program cannot host the tick. ENGAGE: what must an integrator read back before trusting an island? Expected: the firewall registers and the power-sequencing logic. WARN: a block diagram shows what the silicon could do; the board support package decides what is unlocked. LINK: Lecture 2 owns shared-rail power integrity. FLEX: [CORE]; if short, skip step 2's DRIVE comparison and state only that the module's island is closed.}
\end{frame}

\begin{frame}{What can still delay or disable island (c)?}
  \mlsysactivity{Concept question}{2.5 min}
  \centering
  {\normalsize The tick runs on the on-die safety island, from island memory. Which rows can still \emph{delay} or \emph{disable} it?}\\[0.4cm]
  \begin{tabular}{@{}ll@{}}
    \choice{A}{none: the island has its own crystal and rail} & \choice{B}{only the DRAM controller}\\[0.2cm]
    \answerat{2}{C}{die heat, fabric, board supply, reset} & \choice{D}{every row, as for (a)}
  \end{tabular}\\[0.4cm]
  \uncover<2->{\small Heat and fabric can \textbf{delay}; supply and reset can \textbf{disable}.}
  \note{[2.5 min] [CORE] ENGAGE: one minute alone, vote, one minute with a neighbor, vote again. Step 2: C. The island's private caches and memory remove the miss-handling and DRAM rows, and its crystal and rail remove the governed clock and core rail. Four rows remain, die temperature and package (heat can slow it), the on-chip fabric that carries proposals to it (contention can delay the proposals it checks), the board supply (a brownout can disable it), and the reset chain (a host reset can take it down unless the supervisor is independent). Students who chose A treat a dedicated crystal and rail as a complete fault domain; students who chose D forget that the tick stays in island memory. WARN: delay and disable are different failures; a delay breaks the bound, a disable removes the refusal. Land the segment aloud: independence is argued one shared resource at a time. FLEX: [CORE].}
\end{frame}

% =============================================================================
% SECTION: MEMORY
% =============================================================================
\mlsysnote{[0.5 min] [CORE] Second segment: on a shared die, the first place proposal load meets the tick is the memory system. The numbers here put the tick hypothetically on the application SoC, as the chapter does, to price candidates (a) and (b).}
\mlsysdivider{Why does a mostly idle bus still delay a safety read?}{Workload Placement, \S13.3}
\section{Memory}

\begin{frame}{The proposer's weight stream is the long burst ahead of the read}
  \centering
  \safeimg[width=0.80\textwidth,height=4.5cm,keepaspectratio]{socmem.pdf}\\[0.05cm]
  {\small $\text{wait} = \dfrac{\text{bytes queued ahead}}{\text{peak bandwidth}} + \text{row conflicts}\times(t_{\text{precharge}}+t_{\text{activate}})$}\\[0.05cm]
  {\footnotesize\textcolor{midgray}{One backlog (weight tile, camera, telemetry): \val{backlog_bytes}\,B at \val{bw_gbs}\,GB/s + \val{bank_switches} $\times$ \val{conflict_ns}\,ns = \textbf{\val{backlog_us}\,µs}}}
  \mlsyscredit{Workload Placement, \S13.3.1, equation 1 and Napkin Math 13.1 (LPDDR5 at \val{data_rate} MT/s; backlog illustrative).}
  \note{[2 min] [CORE] LINK: Lecture 3 showed that a learned proposer streams its weights from memory on every evaluation, the memory-wall renewal ceiling. NARRATE: the figure is the application SoC with unified memory: the CPU cluster, the NPU that runs the proposer, and the real-time cores where this segment hypothetically places the tick all reach one LPDDR5 controller through the same fabric and last-level cache. On that SoC, the weight stream is the long burst in the same queue as the safety tick's \val{snapshot_kib} KB sensor read. A memory request is issued, granted by the interconnect arbiter, queued at the controller, and served by a DRAM bank. A bank holds one open row; a read to another row of the same bank must close it (precharge) and open the new one (activate), \val{t_rcd_min_ns} ns each in the standard, rounded up to \val{rcd_cycles} cycles of the \val{t_ck_ns} ns command clock: \val{conflict_ns} ns per conflict. The equation is a FIFO wait, the bytes ahead divided by bandwidth plus the conflicts. One backlog, a \val{tile_kib} KB weight tile cut into \val{tile_tx} transactions, a camera transfer, and a telemetry record, waits \val{backlog_us} us (derived). ENGAGE: the bus is idle most of the time; why does the average not help? Expected: the arbiter serves the bursts ahead of the read; averages hide them. FLEX: [CORE].}
\end{frame}

\begin{frame}{\val{n_wcet} queued backlogs cross the declared bound}
  \centering
  \mlsyssteps[width=0.92\textwidth,height=4.4cm,keepaspectratio]{backlog_step}{5}\\[0.1cm]
  \uncover<4->{\small Between \val{n_wcet} and \val{n_dead} backlogs the tick is \textbf{on time} yet \textbf{refutes} the placement.}
  \mlsyscredit{Workload Placement, Napkin Math 13.1 and figure 2 (FIFO model, full channel interleaving; illustrative backlog). Derived in calcs.py.}
  \note{[3.5 min] [CORE] NARRATE: each row is one tick; hatched is the snapshot read waiting behind queued backlogs, green is the \val{compute_after_us} us of evaluation and staging that follows. Step 1: unloaded, the read takes the \val{sense_us} us sense allowance and the tick ends at \val{unloaded_us} us. Step 2: \val{n_sense} backlogs fit in the \val{sense_us} us allowance. Step 3: the declared bound leaves \val{read_room_wcet} us for the read, so \val{n_wcet} backlogs reach \val{wcet} us. Step 4: the deadline less the \val{b_plus_i} us reserve leaves \val{read_room_dead} us for the read: \val{n_dead} backlogs, ending at \val{end_dead} us. Between \val{n_wcet} and \val{n_dead} the tick meets its deadline and still exceeds the WCET, which refutes the placement's premise that proposal load adds nothing. Step 5: if the backlog and the snapshot share one of the \val{channels} channels, each backlog takes \val{backlog1_us} us and \val{n1_wcet} cross the bound. ENGAGE: is plain queueing at full bandwidth the danger? Expected: no; \val{n_sense} backlogs before the sense allowance even breaks. The address mapping is the risk. WARN: an FR-FCFS controller serves row hits first and can keep passing the safety read; this FIFO model is optimistic. FLEX: [CORE]; if short, skip step 5's single-channel case.}
\end{frame}

\begin{frame}{Measured: partitioned victims still slow down up to \val{bechtel_max}$\times$}
  \centering
  \safeimg[width=0.92\textwidth,height=4.3cm,keepaspectratio]{measured.pdf}\\[0.1cm]
  {\small Each victim ran alone on its core with its own cache partition; co-runners streamed to DRAM.}
  \mlsyscredit{Valsan, Yun, Farshchi, RTAS 2016, Fig.~2(b) and 2(e), three co-runners. Bechtel and Yun, RTAS 2019, \S IV-C (Raspberry Pi 3).}
  \note{[2.5 min] [OPTIONAL] NARRATE: published measurements on real multicore chips. A non-blocking cache tracks each outstanding miss in a miss status holding register (MSHR) and parks dirty lines in a writeback buffer; both are small and shared by every core behind the cache. Valsan and colleagues gave the victim its own core and its own cache partition, so its lines were never evicted, and still measured slowdowns of \val{valsan_a9} times on a Cortex-A9, \val{valsan_neh} on Nehalem, and \val{valsan_a15} on a Cortex-A15 when three co-runners streamed reads, and \val{valsan_a15_w} times when they streamed writes. Bechtel and Yun attacked the writeback buffer and slowed a victim on an in-order Cortex-A53 by up to \val{bechtel_max} times. ENGAGE: why does cache partitioning not help? Expected: it divides the lines and leaves the MSHRs and the writeback buffer shared. WARN: these victims were mostly memory stall; the numbers cannot be charged to our tick directly. The next slide turns them into a sensitivity. FLEX: [OPTIONAL].}
\end{frame}

\begin{frame}{Miss handling breaks the bound before bandwidth does}
  \begin{columns}[c]
    \begin{column}{0.58\textwidth}
      \centering\mlsyssteps[width=\textwidth,height=5.3cm,keepaspectratio]{stallshare_step}{2}
    \end{column}
    \begin{column}{0.40\textwidth}
      \small
      Tick $=\val{unloaded_us}\,[1+\varphi(S-1)]$\,µs\\[0.1cm]
      $\varphi \le \dfrac{\val{wcet}/\val{unloaded_us} - 1}{S - 1}$\\[0.25cm]
      {\footnotesize $\varphi$: share of the tick stalled on shared-cache misses.\\ $S$: slowdown of those stalls.\\ Each outstanding miss holds one of a few shared miss registers (MSHRs).}\\[0.25cm]
      \uncover<2->{\footnotesize At $S=\val{s_mshr}$: \textbf{\textcolor{vred}{\val{phi21}\%}} (deadline: \val{phi21_dead}\%)\\ At $S=\val{s_wb}$: \textbf{\textcolor{vred}{\val{phi300}\%}}}
    \end{column}
  \end{columns}
  \mlsyscredit{Workload Placement, \S13.3.2, equation 2. Slowdowns: Valsan et al.\ 2016; Bechtel and Yun 2019. Derived in calcs.py.}
  \note{[2.5 min] [CORE] NARRATE: a cache tracks each outstanding miss in a miss status holding register (MSHR); a core that finds them all taken stalls, whatever its cache partition. Suppose a fraction phi of the unloaded \val{unloaded_us} us tick waits on shared-cache misses, and interference slows that fraction by S while the rest runs unchanged. The tick takes \val{unloaded_us} times (1 + phi(S - 1)). Keeping it under \val{wcet} us gives the bound on the right. The tick can stretch \val{stretch_wcet} times before the WCET and \val{stretch_dead} times before the deadline. Step 1: the red curve; above it, the declared bound is refuted. Step 2: at the measured \val{s_mshr} times, more than \val{phi21} percent miss stall breaks the WCET (\val{phi21_dead} percent for the deadline); at the writeback attack's \val{s_wb} times, \val{phi300} percent. ENGAGE: what removes the multiplier? Expected: a tick that runs from private memory, candidates (c) and (d), or a regulator that bounds S. WARN: this is a sensitivity and predicts nothing; only a profile of the target tells its real stall share. FLEX: [CORE].}
\end{frame}

\begin{frame}{How much miss stall does a \val{s_ten}$\times$ slowdown allow?}
  \mlsysactivity{Estimate}{2 min}
  \begin{columns}[c]
    \begin{column}{0.42\textwidth}
      \begin{mlsyscard}{mlsysaccent}
        \small\textbf{Given}\\[0.05cm]
        Unloaded tick: \val{unloaded_us}\,µs\\
        Declared WCET: \val{wcet}\,µs\\
        Stall slowdown: $S = \val{s_ten}$
      \end{mlsyscard}
    \end{column}
    \begin{column}{0.54\textwidth}
      \small
      What share of the tick may be miss stall?\\[0.3cm]
      \uncover<2->{$\varphi \le \dfrac{\val{wcet}/\val{unloaded_us} - 1}{\val{s_ten} - 1} = \mathbf{\val{phi10}\%}$\\[0.15cm]
      about \textbf{\val{phi10_us}\,µs} of the \val{unloaded_us}\,µs tick}\\[0.2cm]
      \uncover<2->{\footnotesize\textcolor{midgray}{Halving the slowdown more than doubles the room.}}
    \end{column}
  \end{columns}
  \mlsyscredit{Workload Placement, Checkpoint 13.1. Derived in calcs.py.}
  \note{[2 min] [CORE] ENGAGE: two minutes, alone then with a neighbor. Step 2: \val{wcet}/\val{unloaded_us} is \val{stretch_wcet}; less 1 is \val{stretch_wcet_less1}; divided by \val{s_ten_less1} gives \val{phi10} percent, about \val{phi10_us} us of the \val{unloaded_us} us tick (derived). Students who divide by \val{s_ten} instead of \val{s_ten_less1} forget that the unslowed stall is already in the \val{unloaded_us}. Compared with \val{phi21} percent at \val{s_mshr} times, the tolerable share falls as 1/(S - 1), which is why bounding S matters more than shaving stalls. LINK: the remedies segment shows a regulator that measured S near \val{bechtel_w100_1}. Land the segment aloud: arbiters serve the bursts queued ahead of a read, and miss handling in the shared cache threatens the declared bound long before bandwidth does. Rail droop, the third memory-speed coupling, appears on the timescale slide. FLEX: [CORE].}
\end{frame}

% =============================================================================
% SECTION: HEAT
% =============================================================================
\mlsysnote{[0.5 min] [CORE] Third segment: droop and queueing end within microseconds of a burst. Heat does not. Lecture 2 owns thermal duty limits; this segment owns throttling as a placement coupling.}
\mlsysdivider{Why does the check stay late after inference stops?}{Workload Placement, \S13.4}
\section{Heat}

\begin{frame}{Every chip on the board sheds heat into one plate}
  \begin{columns}[c]
    \begin{column}{0.52\textwidth}
      \centering\safeimg[width=\textwidth,height=4.6cm,keepaspectratio]{tesla-cooling-plate.jpg}\\[0.05cm]
      {\footnotesize Tesla Autopilot cooling plate: thermal paste and pads from each chip.}
    \end{column}
    \begin{column}{0.46\textwidth}
      {\footnotesize\textbf{Jetson AGX Orin, Linux device tree:}}\par\vspace{0.08cm}
      \begin{logblock}[7.5]
tj-thermal \{\\
\ind{2}polling-delay = <1000>;\\
\ind{2}trips \{\\
\ind{4}tj\_trip\_active0: active-0 \{\\
\ind{6}temperature = <75000>;\\
\ind{6}hysteresis = <4000>;\\
\ind{4}\};\\
\ind{4}tj\_trip\_active1: active-1 \{\\
\ind{6}temperature = <95000>;\\
\ind{6}hysteresis = <4000>;\\
\ind{4}\};
      \end{logblock}
      \vspace{0.1cm}
      {\footnotesize One junction zone for the die, polled every \val{orin_poll_ms}\,ms.}
    \end{column}
  \end{columns}
  \mlsyscredit{Photo: Steve Jurvetson, CC BY 2.0, Wikimedia Commons. Code: Linux v6.12, arch/arm64/boot/dts/nvidia/tegra234-p3701-0000.dtsi (excerpt).}
  \note{[2 min] [OPTIONAL] NARRATE: left, the liquid-cooled plate from a Tesla Autopilot computer; each square of paste and pad marks a chip that dumps its heat into the same metal. Right, real source from the upstream Linux device tree for the Jetson AGX Orin module. One thermal zone, tj, the junction, covers the die; it is polled once a second, a thousand permission ticks; its trips at 75 and 95\,°C drive the fan with 4\,°C of hysteresis. Temperatures are in thousandths of a degree. ENGAGE: why does a governor need hysteresis? Expected: without a gap between trip and clear, the fan or clock would chatter around the threshold. WARN: these trips drive the fan (active trips); our example's clock throttle is a different trip; the point is that one sensor and one slow loop serve every block on the die. FLEX: [OPTIONAL].}
\end{frame}

\begin{frame}{A halved clock stays halved for \val{recover_s}\,s after the load stops}
  \centering
  \mlsyssteps[width=0.92\textwidth,height=5.05cm,keepaspectratio]{hangover_step}{2}\\[0.02cm]
  {\small Thermal throttling hangover: the governor follows temperature, which lags the load by seconds.}
  \mlsyscredit{Workload Placement, Napkin Math 13.2 and figure 3 (illustrative: $R=\val{theta}$\,K/W, $\tau=\val{tau_s}$\,s, \val{p_quiet}\,W idle, \val{p_load}\,W loaded, \val{t_amb}\,°C bay). Derived in calcs.py.}
  \note{[3.5 min] [CORE] NARRATE: the worked example behind the hook. Junction steady states: \val{t_amb} + \val{p_quiet} W times \val{theta} K/W = \val{t_quiet}\,°C idle; \val{t_amb} + \val{p_load} times \val{theta} = \val{t_loaded}\,°C with inference, above the \val{t_trip_c}\,°C trip, reached after \val{t_trip_s} s. Top panel: after inference stops, the junction cools exponentially with the \val{tau_s} s time constant toward \val{t_quiet}; the clock returns only at the \val{t_clear_c}\,°C clear point, after \val{tau_s} ln(\val{trip_head}/\val{clear_head}) = \val{recover_s} s, \val{recover_ticks} ticks. Step 1: the nominal tick, \val{tick_cycles} cycles at \val{f_nom_ghz} GHz, \val{wcet} us, plus \val{b_plus_i} us, ends at \val{resp_nom}. Step 2: at \val{f_thr_mhz} MHz the same code takes \val{thr_us} us and ends at \val{resp_thr}, \val{past_us} us past the deadline. ENGAGE: shedding proposal load is the remedy for memory contention; why does it fail here? Expected: the clock follows temperature, which lags by seconds. WARN: each late refusal slips about one tick, which the stopping budget could absorb; the placement still fails, because the loaded tick has crossed the declared bound and nothing limits a real governor to halving. The lever is the heat path. Any $R$ below \val{theta_max} K/W never trips (derived), or an independently clocked placement. FLEX: [CORE]; if short, give only step 2's \val{resp_thr} us end time.}
\end{frame}

\begin{frame}{Couplings act from nanoseconds to seconds around a 1\,ms tick}
  \centering
  \safeimg[width=0.92\textwidth,height=4.6cm,keepaspectratio]{ladder.pdf}\\[0.15cm]
  {\small Rail droop: \val{di_a}\,A in \val{rise_ns}\,ns $\rightarrow$ \textbf{\val{droop_mv}\,mV} on a \val{v_nom}\,V rail, gate delay \textbf{+\val{delay_pct}\%} (modeled).}
  \mlsyscredit{Workload Placement, \S13.4 ladder and \S13.3.3 (illustrative package: \val{l_ph}\,pH, \val{r_mohm}\,m$\Omega$; alpha-power model, Sakurai and Newton 1990).}
  \note{[2 min] [OPTIONAL] NARRATE: the couplings on one log axis. Rail droop rises in about \val{rise_ns} ns. A \val{di_a} A step through \val{l_ph} pH of package inductance gives \val{droop_l_mv} mV and through \val{r_mohm} milliohms \val{droop_r_mv} mV, \val{droop_mv} mV from \val{v_nom} V; the alpha-power model stretches gate delay \val{delay_pct} percent while it lasts. A bank conflict costs \val{conflict_ns} ns, a backlog \val{ladder_backlog_us} us, the protected tick \val{t_tick_ms} ms, thermal recovery \val{recover_s} s, and chassis heat accumulates over hours. ENGAGE: a latency spike correlates with rail voltage; another with temperature. Which mechanism is each? Expected: power delivery, then the governor. WARN: the droop lasts nanoseconds, so it does not lengthen the tick \val{delay_pct} percent; it threatens setup time on paths switching during the droop, and many processors stretch the clock to survive it. LINK: Lecture 2 owns shared-rail power integrity. Land the segment: heat couples the paths on a timescale of seconds. FLEX: [OPTIONAL].}
\end{frame}

% =============================================================================
% SECTION: REMEDIES
% =============================================================================
\mlsysnote{[0.5 min] [CORE] Fourth segment: a loaded test caught a coupling. What fixes it, at what price, and once the paths are separated, how do they talk?}
\mlsysdivider{What does each fix cost, and how do separated paths talk?}{Workload Placement, \S13.5 and \S13.6}
\section{Remedies}

\begin{frame}{Each remedy fixes one coupling and costs a different currency}
  \begin{columns}[c]
    \begin{column}{0.74\textwidth}
      \footnotesize
      \begin{tabular}{@{}>{\raggedright\arraybackslash}p{2.3cm}>{\raggedright\arraybackslash}p{2.4cm}>{\raggedright\arraybackslash}p{2.2cm}>{\raggedright\arraybackslash}p{2.4cm}@{}}\toprule
        \textbf{Remedy} & \textbf{Fixes} & \textbf{Pays in} & \textbf{Leaves}\\\midrule
        Reservation, bandwidth caps & CPU time, memory demand & idle capacity & cache structures, droop\\\addlinespace[0.12cm]
        Cache coloring, bank partition & working-set thrash & smaller working sets & controller queue, MSHRs\\\addlinespace[0.12cm]
        Workload reduction & peak traffic and heat & policy fidelity & clock, reset, current steps\\\addlinespace[0.12cm]
        \textbf{Relocation} & shared rail, clock, reset & board area, a link & harness, battery\\\bottomrule
      \end{tabular}
    \end{column}
    \begin{column}{0.24\textwidth}
      \centering\safeimg[width=\textwidth,height=2.6cm,keepaspectratio]{tesla-hw25-tricore.jpg}\\[0.05cm]
      {\footnotesize The relocated MCU on a real board}
    \end{column}
  \end{columns}
  \vspace{0.1cm}
  {\small A remedy that misses the coupling strands capacity and keeps the overrun.}\\[0.1cm]
  \mlsyscredit{Workload Placement, \S13.6, table 3 (condensed). Photo: Steve Jurvetson, CC BY 2.0, Wikimedia Commons.}
  \note{[2.5 min] [CORE] NARRATE: four remedies, each matched to a coupling. Reservation gives the tick fixed time slots and caps each core's memory traffic; the price is capacity withheld even when the enforcer is idle. Coloring and bank partitioning give each domain private cache sets and DRAM banks through the page allocator; the price is a smaller working set, and the MSHRs and controller queue stay shared. Workload reduction shrinks traffic and heat by quantizing the model, lowering the policy rate, or pruning logs; the price is fidelity, and Lecture 3 showed a quantized model is a new model that needs fresh testing. Keeping the \val{lease_ms} ms lease alive needs admissions faster than \val{min_rate_hz} Hz. Relocation moves the tick to a separate MCU, the TriCore in the photo; the price is board area and a link. ENGAGE: the stall lives in the DRAM queue; what does reserving a whole CPU core buy? Expected: nothing for the stall. FLEX: [CORE].}
\end{frame}

\begin{frame}{Predict: which fix restores a victim slowed \val{bechtel_unreg}$\times$?}
  \mlsysactivity{Predict}{2 min}
  \centering
  \mlsyssteps[width=0.92\textwidth,height=3.9cm,keepaspectratio]{regulation_step}{2}\\[0.1cm]
  {\small The victim already has its own core and its own cache partition.}\\[0.2cm]
  {\normalsize\choice{A}{a faster core}\qquad\choice{B}{a larger partition}\qquad\answerat{2}{C}{cap the attackers' writes}}
  \mlsyscredit{Bechtel and Yun, RTAS 2019, \S VI (MemGuard extended to writebacks; Raspberry Pi 3, three BwWrite attackers).}
  \note{[2 min] [CORE] NARRATE: Step 1: the same attack as the memory segment, in the paper's regulation experiment. A victim on its own core and cache partition slows \val{bechtel_unreg} times when three cores stream writes; the \val{bechtel_max} times on the measured slide is the paper's largest slowdown across its benchmarks (Section IV-C), and \val{bechtel_unreg} times is the unregulated baseline of this experiment (Section VI). ENGAGE: one minute, vote. Step 2: C. Bechtel and Yun extended per-core bandwidth regulation, the MemGuard idea of Yun and colleagues, to count writebacks. Each attacker core gets a budget per period, performance counters count its traffic, and an interrupt stalls it when the budget runs out. With writes capped at 100 MB/s the victim slowed only \val{bechtel_w100} times; at 50 MB/s, \val{bechtel_w50} times. Students who chose B think in capacity; the partition was already there. A attacks the wrong resource. WARN: regulation bounds the multiplier S of the previous segment only if both reads and writebacks are charged, and the cost is throughput withheld from the proposer even when the robot is parked. FLEX: [CORE].}
\end{frame}

\begin{frame}[fragile]{A static-partitioning hypervisor binds cores and memory at boot}
  \begin{columns}[c]
    \begin{column}{0.30\textwidth}
      {\footnotesize\textbf{\textcolor{vnavy}{VM 1: Linux}}}\par\vspace{0.05cm}
      \begin{logblock}[8]
.platform = \{\\
\ind{2}.cpu\_num = 3,\\
\ind{2}.region\_num = 1,\\
\ind{2}[...]\\
\ind{4}.base = 0x00000000,\\
\ind{4}.size = 0x40000000
      \end{logblock}
      {\footnotesize \val{bao_linux_cpus} cores, \val{bao_linux_gib}\,GiB}
    \end{column}
    \begin{column}{0.30\textwidth}
      {\footnotesize\textbf{\textcolor{vgreen}{VM 2: FreeRTOS}}}\par\vspace{0.05cm}
      \begin{logblock}[8]
.platform = \{\\
\ind{2}.cpu\_num = 1,\\
\ind{2}.region\_num = 1,\\
\ind{2}[...]\\
\ind{4}.base = 0x0,\\
\ind{4}.size = 0x8000000
      \end{logblock}
      {\footnotesize \val{bao_rtos_cpus} core, \val{bao_rtos_mib}\,MiB}
    \end{column}
    \begin{column}{0.36\textwidth}
      \small Shared: one \val{bao_shmem_kib}\,KiB region for messages.\\[0.25cm]
      \uncover<2->{\small\textcolor{vred}{Not named in the file:} cache, DRAM controller, rail, clock.}
    \end{column}
  \end{columns}
  \mlsyscredit{Bao hypervisor demos, demos/linux+freertos/configs/zcu104.c (excerpt; github.com/bao-project/bao-demos). Martins et al., NG-RES 2020.}
  \note{[2 min] [OPTIONAL] NARRATE: a real configuration file for Bao, a static-partitioning hypervisor, on a Xilinx ZCU104 board. At boot it binds \val{bao_linux_cpus} cores and \val{bao_linux_gib} GiB to a Linux guest and \val{bao_rtos_cpus} core and \val{bao_rtos_mib} MiB to a FreeRTOS guest, with one \val{bao_shmem_kib} KiB shared region and its interrupt for messages; there is no scheduler at run time. This is the standard way to host a real-time enclave beside Linux on one die, and it matches mixed-criticality scheduling, in which the high-criticality enforcer keeps its budget, and the low-criticality proposer forfeits lease renewal when it runs late. Step 2: read what the file does not mention. ENGAGE: which interference rows from the inventory does this file close? Expected: cores and address spaces only; Bao can also color caches, but this demo does not, and nothing in it bounds the DRAM queue, the rail, or the clock. WARN: a partitioned configuration is a spatial claim; the temporal claim still needs a loaded test. FLEX: [OPTIONAL].}
\end{frame}

\begin{frame}{The board link costs \val{link_ms}\,ms of a \val{d2a_ms}\,ms allocation}
  \centering
  \mlsyssteps[width=0.94\textwidth,height=4.1cm,keepaspectratio]{boardlink_step}{2}\\[0.1cm]
  {\small Chunk \val{chunk_bytes}\,B on \val{spi_mhz}\,MHz SPI.\enspace At \val{v_aisle}\,m/s, \val{spare_mm}\,mm of spare pays for at most \val{ceil_ms}\,ms.}\\[0.1cm]
  \uncover<2->{\footnotesize Off board: RT-2's 5B model at about \val{rt2_5b_hz}\,Hz returns a plan every \val{rt2_5b_ms}\,ms, \val{rt2_ratio}$\times$ the \val{renew_window_ms}\,ms renewal window.}
  \mlsyscredit{Workload Placement, Napkin Math 13.3 (illustrative allowances) and \S13.6.3 (Brohan et al.\ 2023). Spare: Safety Enforcement, Napkin Math 12.1. Derived in calcs.py.}
  \note{[2.5 min] [CORE] LINK: Lecture 12 left \val{spare_mm} mm of spare clearance at the \val{v_aisle} m/s aisle speed; Lecture 8's latency waterfall left one interval unallocated, from perception dispatch to admission. NARRATE: Step 1: the allocation. Inference \val{chunk_infer_ms} ms (charged in full, conservatively); the link, 8 times \val{chunk_bytes} bytes at \val{spi_mhz} MHz, \val{link_ms} ms; deserialization \val{deser_ms}, validation \val{valid_ms}, activation \val{activate_ms}, and up to one tick of waiting: \val{d2a_ms} ms, during which the base travels \val{d2a_mm} mm on aging evidence. Step 2: \val{spare_mm} mm divided by \val{v_aisle} m/s gives a \val{ceil_ms} ms ceiling, so \val{reserve_ms} ms, or \val{reserve_mm} mm, remains. The link that buys separation is about \val{link_share} percent; inference and validation dominate. Another inference's worth, \val{plus_infer_ms} ms, would exhaust the spare. ENGAGE: why can a data-center model not renew the lease? Expected: RT-2 ran at \val{rt2_55b_low} to \val{rt2_55b_high} Hz for 55B and about \val{rt2_5b_hz} Hz for 5B over the network; \val{rt2_5b_ms} ms is \val{rt2_ratio} times the \val{renew_window_ms} ms window. WARN: jitter spends the \val{jitter_ms} ms renewal margin of Lecture 4's lease; a \val{stretch_pct} percent inference stretch uses all of it. FLEX: [CORE]; if short, skip step 2's RT-2 comparison.}
\end{frame}

\begin{frame}{Dispatch to admission leaves \val{reserve_mm} mm of the stopping budget}
  \centering
  \mlsyskitfig[width=0.95\textwidth,height=4.8cm,keepaspectratio]{budget_bar_aisle_step8}\\[0.05cm]
  {\small \val{d2a_ms}~ms at \val{v_aisle}~m/s is \val{d2a_mm}~mm, leaving a \val{reserve_mm}~mm reserve, or \val{reserve_ms}~ms.}
  \mlsyscredit{Workload Placement, Napkin Math 13.3; Safety Enforcement, Napkin Math 12.1. Stopping-budget bar: course kit, MLSysIM values.}
  \note{[0.5 min] [CORE] NARRATE: the stopping-budget bar at the \val{v_aisle} m/s aisle speed, as Lecture 12 left it, gains today's row: dispatch to admission, \val{d2a_ms} ms, \val{d2a_mm} mm of travel on aging evidence. The reserve left is \val{reserve_mm} mm, \val{reserve_ms} ms. LINK: Lecture 14 spends that reserve on takeovers. FLEX: [CORE].}
\end{frame}

\begin{frame}[fragile]{The enforcer reads the mailbox once per tick and never waits}
  \begin{columns}[c]
    \begin{column}{0.60\textwidth}
      {\footnotesize\textbf{Linux reader: retries until a clean copy}}\par\vspace{0.08cm}
      \begin{logblock}[7.5]
\mbox{do \{}\\
\mbox{\ind{8}seq = read\_seqcount\_begin(\&foo\_seqcount);}\\[0.1cm]
\mbox{\ind{8}/* ... [[read-side critical section]] ... */}\\[0.1cm]
\mbox{\} while (read\_seqcount\_retry(\&foo\_seqcount, seq));}
      \end{logblock}
    \end{column}
    \begin{column}{0.37\textwidth}
      \small\textbf{\textcolor{vgreen}{Boundary contract: one attempt per tick}}\\[0.1cm]
      Sequence odd or changed: \textbf{discard} the record this tick.\\[0.1cm]
      The admitted chunk runs on while its lease holds; then the fallback.
    \end{column}
  \end{columns}
  \vspace{0.3cm}
  {\small A retry loop would hand the deadline to a stalled proposer; a windowed watchdog checks that the enforcer itself stays alive.}
  \mlsyscredit{Linux v6.12, Documentation/locking/seqlock.rst (excerpt). Workload Placement, \S13.5.1; The Nervous System, the single-attempt seqlock.}
  \note{[2 min] [CORE] LINK: Lecture 4 owns the freshest-proposal mailbox and its seqlock. NARRATE: left, real code from the Linux kernel documentation. A seqlock writer bumps a sequence number before and after writing; the reader samples it, copies the data, and retries if the number was odd or changed. That loop is correct for Linux, but it can spin as long as a writer keeps writing. Right, the enforcer's rule is one attempt per tick. A torn or in-flight record is discarded; the chunk admitted earlier keeps running only while its own lease, bounds, and state checks hold, and when it expires the fallback ladder of Lecture 12 takes over. ENGAGE: a proposer crashes halfway through publishing. What does the enforcer do on the next tick? Expected: sees an odd sequence, discards, keeps the current command until its lease lapses. WARN: placement adds two obligations the protocol does not cover, namely cache and memory settings that keep the copy coherent, and a largest-payload copy that fits inside the WCET. FLEX: [CORE].}
\end{frame}

\begin{frame}{A windowed watchdog trips on early feeds as well as late ones}
  \centering
  \mlsyssteps[width=0.90\textwidth,height=4.4cm,keepaspectratio]{watchdog_step}{3}\\[0.1cm]
  {\small Early feed: skipped checks or a fast clock.\enspace Late feed: a stalled enforcer.}\\[0.05cm]
  \uncover<3->{\footnotesize Machine's late trip: \val{wd_ms}\,ms, \textbf{\val{wd_mm}\,mm} of travel at \val{v_aisle}\,m/s before a dead enforcer is detected.}
  \mlsyscredit{Workload Placement, \S13.5.2 and equation 5. Barr 2013 (Bookout v.\ Toyota). Derived in calcs.py.}
  \note{[2.5 min] [CORE] NARRATE: Step 1: an ordinary watchdog resets the chip if no feed arrives by a deadline. Lecture 15 tells the Toyota case, in which a watchdog fed from a timer interrupt kept feeding after the control task died. Step 2: a windowed watchdog also rejects feeds that come too early. An early feed means control flow skipped checks, a corrupted program counter jumped to the feed routine, or a clock fault sped the core. Step 3: the window must accept every healthy feed, so it spans the period plus or minus the spread between the earliest and latest tick completion and the period jitter, all measured under load. The machine keeps a \val{wd_ms} ms late trip, several ticks past the window, which is \val{wd_mm} mm of base travel at \val{v_aisle} m/s (derived). ENGAGE: why not use the unloaded \val{unloaded_us} us as the earliest completion? Expected: early-exit paths finish faster. WARN: a valid feed proves timing only; it cannot show that every check ran; challenge-response checks and an external supervisor cover that. If the MCU dies it cannot command a powered stop; the drives' certified stop must work without it. LINK: Lecture 4's self-watchdog had only the late trip. FLEX: [CORE].}
\end{frame}

% =============================================================================
% SECTION: REFUTATION
% =============================================================================
\mlsysnote{[0.5 min] [CORE] Land the fourth segment aloud: each remedy matches one coupling, and relocation costs a board link the budget can afford. Last segment: what measurement could prove the chosen placement wrong?}
\mlsysdivider{What measurement could prove the placement wrong?}{Workload Placement, \S13.7 and \S13.8}
\section{Refutation}

\mlsysnote{[1.5 min] [CORE] NARRATE: the key relation of the lecture. Loaded execution time is isolated execution time plus the delay from hardware shared with the proposer. The isolated time cannot exceed the declared WCET by definition, and the placement claims the shared term is zero. So the declared WCET is the falsification threshold, and the deadline is not. A loaded tick of \val{ex_tick} us refutes the placement even though it ends \val{ex_before_deadline} us before the \val{deadline} us deadline, because every margin downstream, down to the distance budgeted per tick, was computed on the premise it just broke. ENGAGE: does a loaded tick below \val{wcet} prove the shared term is zero? Expected: no; a test samples only the states it reached. WARN: a single tick at \val{v_aisle} m/s is \val{tick_mm} mm, which tempts students to shrug. The threat is an open-ended stall, which no bound limits once the premise has failed. FLEX: [CORE].}
\mlsysfocus{The WCET is the threshold}{%
  $\Tload = \Tiso + \Dint,\qquad \Tiso \le \Twcet < \Tdead$\\[0.5cm]
  {\large The declared WCET is the falsification threshold.\\ A loaded tick above it refutes the placement,\\ inside the deadline or not.}\\[0.8cm]\mbox{}%
}

\begin{frame}{Which claim does a \val{cp_us}\,µs loaded cycle refute?}
  \mlsysactivity{Concept question}{2.5 min}
  \centering
  {\normalsize Shared-die candidate, loaded campaign: one cycle of \textbf{\val{cp_us}\,µs}.\\ Declared WCET \val{wcet}\,µs; deadline \val{deadline}\,µs; \val{b_plus_i}\,µs of the MCU's own tasks.}\\[0.35cm]
  \begin{tabular}{@{}l@{}}
    \choice{A}{none: the response still meets the deadline}\\[0.1cm]
    \answerat{2}{B}{the claim that shared hardware adds no delay}\\[0.1cm]
    \choice{C}{the deadline itself}\\[0.1cm]
    \choice{D}{only the unloaded \val{unloaded_us}\,µs estimate}
  \end{tabular}\\[0.3cm]
  \uncover<2->{\small \val{cp_us} $-$ \val{wcet} = \textbf{\val{cp_over}\,µs} over; \val{cp_us} + \val{b_plus_i} = \val{cp_resp}\,µs, \textbf{\val{cp_left}\,µs} inside the deadline.}
  \mlsyscredit{Workload Placement, Checkpoint 13.3. Derived in calcs.py.}
  \note{[2.5 min] [CORE] ENGAGE: vote, discuss, vote. Step 2: B. The cycle exceeds the declared bound by \val{cp_over} us; with the \val{b_plus_i} us of local blocking and interference the response is \val{cp_resp} us, \val{cp_left} us inside the deadline (derived). The deadline held this time, but the joint claim that the isolated baseline is bounded and proposal load adds nothing is false. Students who chose A judge by the deadline; that is the error this lecture exists to prevent. D is half right, because an overrun alone cannot tell an underestimated baseline from interference; attributing it needs the baseline justified separately. Either way the placement's timing premise fails. WARN: only \val{cp_left} us remained; the next stall is not bounded by anything. FLEX: [CORE].}
\end{frame}

\begin{frame}{A loaded trace can refute a placement; a model cannot prove one}
  \begin{columns}[c]
    \begin{column}{0.56\textwidth}
      \centering\mlsyssteps[width=\textwidth,height=5.0cm,keepaspectratio]{trace_step}{2}\\[0.02cm]
      {\footnotesize\textcolor{midgray}{Illustrative $10^{\val{trace_exp}}$-tick trace: \val{trace_h}\,h \val{trace_min}\,min \val{trace_sec}\,s.}}
    \end{column}
    \begin{column}{0.42\textwidth}
      \begin{mlsyscard}{vnavy}\footnotesize
        ``The WCET of a software component or task may increase significantly when other software components or tasks are executing in parallel on the other cores of an MCP.''\\[0.05cm]
        \hfill\textcolor{midgray}{FAA AC 20-193, \S5.5.3}
      \end{mlsyscard}
      \vspace{0.1cm}
      {\footnotesize WCET from static analysis, measurement plus margin, hybrid, or extreme-value fits; local worst cases need not compose.}
    \end{column}
  \end{columns}
  \mlsyscredit{Workload Placement, \S13.7.2 and \S13.7.3 (percentiles illustrative). FAA, Use of Multi-Core Processors, AC 20-193, 1/8/24.}
  \note{[2.5 min] [OPTIONAL] NARRATE: Step 1: cold, idle proposer on the rejected shared-die placement: median \val{cold_p50} us, p99 \val{cold_p99}, maximum \val{cold_max}, all under \val{wcet}. Step 2: warm and contended: median \val{hot_p50}, p99 \val{hot_p99}, p99.9 \val{hot_p999}, maximum \val{hot_max}. The p99 alone exceeds the declared bound, and with \val{b_plus_i} us charged it passes the deadline. The trace reaches the verdict the thermal model reached, without relying on the model. The percentiles are illustrative. Read the FAA sentence. Aviation's multicore guidance (2024) requires identifying every interference channel and verifying its mitigation on the target in the final configuration. The four WCET methods are listed; on our in-order MCU with private memory, static analysis is tractable; on an application multicore, timing anomalies mean a local speedup can lengthen the whole run. ENGAGE: the campaign must trigger the refusal branch under full load; why? Expected: the cold refusal path runs rarely and may fetch code from DRAM exactly when a diagnostic dump floods the controller. WARN: zero violations in \val{trace_millions} million ticks supports the bound only for the states tested. FLEX: [OPTIONAL]; cut first when short, since the predict before it already reaches the verdict.}
\end{frame}

\begin{frame}{The placement record's loaded field is still empty}
  \centering\footnotesize
  \begin{tabular}{@{}ll@{}}\toprule
    \textbf{Field} & \textbf{Running machine}\\\midrule
    execution\_domain\_id & dedicated dual-core lockstep MCU (pending)\\
    nominal\_period\_ns / hard\_deadline\_ns & \val{tick_ns} / \val{deadline_ns}\\
    unloaded\_observed\_max\_ns & \texttt{0xFFFFFFFFFFFFFFFF} (estimate \val{unloaded_us}\,µs in report)\\
    \textbf{loaded\_observed\_max\_ns} & \textbf{\textcolor{vred}{\texttt{0xFFFFFFFFFFFFFFFF}}} (no campaign yet)\\
    shared\_axi\_channel\_qos & TCM-local bus; bounds pending\\
    max\_pdn\_droop\_thermal\_trip & pending device characterization\\
    enforcement\_record\_digest & SHA-256 of the enforcement record\\
    test\_report\_digest & unset until a report exists\\\bottomrule
  \end{tabular}\\[0.25cm]
  {\small Pass rule, fixed before testing: loaded cycles $\le$ \val{wcet}\,µs; brake onset $\le$ \val{renewal_onset_ms}\,ms.}
  \mlsyscredit{Workload Placement, \S13.8, table 4 (selected fields).}
  \note{[2.5 min] [CORE] NARRATE: the record this lecture adds to the chain. It consumes Lecture 12's enforcement record, taking its declared WCET and deadline, and binds them to evidence: the execution domain, the period and deadline, the unloaded and loaded maxima, the bus arbitration settings, the droop and thermal limits, the mitigations, and digests of the enforcement record and the test report. On the running machine the measurement fields hold all ones, a sentinel that blocks release sign-off. The pass criterion is written before the campaign. Every loaded cycle must stay at or below \val{wcet} us, and powered deceleration must start within \val{renewal_onset_ms} ms of the last valid renewal. The external report lists what was not tested, such as bays above \val{t_amb}\,°C, multi-bit errors in tightly coupled memory, and silicon aging. ENGAGE: why fix the criterion before testing? Expected: so the test can fail. LINK: Lecture 15 turns these arbitration points and rails into a fault list; Lecture 16 shows that certification authorizes one configuration. FLEX: [CORE].}
\end{frame}

% =============================================================================
% CLOSE
% =============================================================================
\begin{frame}{Why the check stayed late on its own core}
  \mlsysactivity{Retrieval}{2 min}
  \begin{columns}[c]
    \begin{column}{0.36\textwidth}
      \centering\safeimg[width=\textwidth,height=2.2cm,keepaspectratio]{thumb_hook_step3.pdf}\\[0.05cm]
      {\footnotesize \val{resp_thr}\,µs ticks for \val{recover_s}\,s after the stop}
    \end{column}
    \begin{column}{0.62\textwidth}
      \only<1|handout:0>{\normalsize Close your notes. Name three couplings the board diagram hid, and the remedy for each.}%
      \only<2->{\footnotesize
      \begin{tabular}{@{}>{\raggedright\arraybackslash}p{4.0cm}>{\raggedright\arraybackslash}p{4.6cm}@{}}\toprule
        \textbf{Coupling} & \textbf{Remedy}\\\midrule
        Heat, shared governor & heat path; own clock\\
        Cache misses, DRAM queue & regulation; private memory\\
        Rail, reset & isolated rail; supervisor\\\bottomrule
      \end{tabular}\\[0.15cm]}
      \uncover<3->{\small The discrete MCU removes all three; the battery and the link remain, and the loaded test has not run.}
    \end{column}
  \end{columns}
  \note{[2.5 min] [CORE] Back to the hook. Step 1: retrieval, 60 seconds. Step 2: the governor followed temperature, which lags the load, so the clock stayed halved for \val{recover_s} s; the memory queue and the shared-cache miss structures sat beneath every privilege check; and a shared rail and reset chain could silence both paths at once. Each has a remedy priced in a different currency. Step 3: the running machine chose the discrete MCU (d), because the module would not expose its island. That removes the host governor, the shared DRAM, and the shared rail, keeps the battery and adds a \val{link_ms} ms link the budget can afford, and still owes the loaded campaign whose result the placement record does not yet hold. ENGAGE: what one measurement would most change your trust in this machine? Expected: the loaded maximum under the factorial campaign. FLEX: [CORE].}
\end{frame}

\begin{frame}{Takeaways}
  \thumbline{thumb_board_step4.pdf}{Independence is argued one shared resource at a time; a separate core or chip is a hypothesis.}\\[0.12cm]
  \thumbline{thumb_stallshare_step2.pdf}{Arbiters serve the bursts ahead of a read; at a \val{s_mshr}$\times$ slowdown, \val{phi21}\% miss stall breaks the \val{wcet}\,µs bound.}\\[0.12cm]
  \thumbline{thumb_hangover_step2.pdf}{Heat holds a halved clock for \val{recover_s}\,s after the load stops; shedding load cannot rescue the tick.}\\[0.12cm]
  \thumbline{thumb_boardlink_step2.pdf}{Each remedy has a price; relocation's link costs \val{link_ms}\,ms of a \val{d2a_ms}\,ms allocation under \val{ceil_ms}\,ms.}\\[0.12cm]
  \thumbline{thumb_trace_step2.pdf}{A loaded tick above the declared WCET refutes the placement; the placement record carries what was tested.}
  \note{[1 min] [CORE] One line per segment. Ask students which slide taught each. Close Part III aloud. Proposal is not permission only while the proposer cannot reach the permission decision through anything the two share, and on a real machine that includes memory, supply, clock, reset, and heat.}
\end{frame}

\begin{frame}{Open questions}
  \begin{columns}[c]
    \begin{column}{0.62\textwidth}
      \small
      \begin{itemize}\setlength{\itemsep}{0.12cm}
        \item Should a safety argument rest on a silicon feature a vendor may withhold on the next module?
        \item How many hours of loaded testing should a release require, and who decides?
        \item Which proposals, if any, may a robot accept from a data center?
      \end{itemize}
    \end{column}
    \begin{column}{0.36\textwidth}
      \centering\safeimg[width=\textwidth,height=3.3cm,keepaspectratio]{tesla-fsd-board.jpg}\\[0.05cm]
      {\footnotesize Both SoCs draw on one board supply.}
    \end{column}
  \end{columns}
  \mlsyscredit{Photo: Votpuske, CC BY 4.0, Wikimedia Commons.}
  \note{[0.5 min] [OPTIONAL] No single answer. The first follows from table 2: the Jetson Thor module does not support its island while IGX Thor does. The second: \val{trace_millions} million ticks is under three hours on one machine, and it bounds nothing across a fleet. The third: FogROS2 cut grasp planning from \val{fogros_local_s} s to \val{fogros_off_s} s by offloading, a gain for second-scale planning but useless for a \val{lease_ms} ms lease. Assign one as a short written response. FLEX: [OPTIONAL].}
\end{frame}

\begin{frame}{Read and explore}
  \small
  \begin{itemize}\setlength{\itemsep}{0.02cm}
    \item \textbf{Read:} Workload Placement, \S13.1 to \S13.10; Checkpoints 13.1 to 13.3; Napkin Math 13.1 to 13.3 by hand.
    \item Rushby, \emph{Partitioning in Avionics Architectures}, NASA, 1999: spatial and temporal partitioning.
    \item Yun et al., ``MemGuard,'' RTAS 2013: per-core memory-bandwidth regulation.
    \item Valsan, Yun, Farshchi, ``Taming Non-blocking Caches,'' RTAS 2016: the MSHR measurements.
    \item Wilhelm et al., ``The Worst-Case Execution-Time Problem,'' ACM TECS 2008.
    \item FAA, AC 20-193, \emph{Use of Multi-Core Processors}, 2024.
  \end{itemize}
  \note{[0.5 min] [CORE] The chapter first, with its checkpoints. Rushby separates the spatial and temporal obligations the inventory checks. MemGuard is the regulator behind the predict slide. Valsan and colleagues measured the \val{s_mshr} times. Wilhelm and colleagues survey WCET methods. AC 20-193 is what a regulator asks of a shared-die timing bound. FLEX: add a lab that measures a co-runner's effect on a timed loop on any multicore board.}
\end{frame}

\begin{frame}[c]
  \begin{columns}[c]
    \begin{column}{0.34\textwidth}
      \centering\safeimg[width=\textwidth,height=4.2cm,keepaspectratio]{thumb_next_stack_governance.pdf}
    \end{column}
    \begin{column}{0.62\textwidth}
      {\small\textcolor{lightgray}{NEXT LECTURE \textbar{} SUPERVISORY INTERVENTION}}\\[0.5cm]
      {\LARGE\bfseries Why can handing a moving robot to a person cause the collision the handover was meant to prevent?}
    \end{column}
  \end{columns}
  \note{[0.5 min] [CORE] Every delay today belonged to the machine and could be bounded or refuted on the chassis. Lecture 14 opens Part IV with a delay no silicon placement can bound, a human taking over, spending the same stopping budget, as the takeover delay in the stopping budget.}
\end{frame}

\end{document}
